\section{Introduction}\label{sec:introduction}

For a convex body $K\subset\R^n$, its lattice covering density is
\[
 \theta_L(K)=\inf\left\{
 \frac{\vol_n(K)}{\det\Lambda}:\Lambda\text{ is a full-rank lattice and }
 K+\Lambda=\R^n\right\}.
\]
Here a convex body is compact with nonempty interior, and
$\det\Lambda$ denotes lattice covolume. The density is the average number
of translates covering a point in a fundamental cell: integrating the
covering multiplicity over that cell gives $\vol_n(K)$. Thus density
measures the overlap needed to cover space when the centers must form
one lattice. We prove the following bound.

\begin{theorem}\label{thm:main}
There is an absolute constant $C>0$ such that, for every integer
$n\ge2$ and every convex body $K\subset\R^n$, one full-rank lattice
$\Lambda$ satisfies
\[
 K+\Lambda=\R^n,\qquad
 \frac{\vol_n(K)}{\det\Lambda}\le Cn\log n.
\]
No symmetry or boundary regularity is assumed.
\end{theorem}

\paragraph{History and significance.}
Rogers proved that every convex body in dimension $n\ge3$ admits a
translative covering of density at most
$n\log n+n\log\log n+5n$~\cite{Rog57}.
The centers in that construction may be a finite union of lattice
cosets. Fejes T\'oth later obtained the same leading density using only
$O(\log n)$ cosets of one lattice~\cite{FT09}.
Requiring the centers to form one lattice is a stronger
constraint, and a translative bound does not establish
Theorem~\ref{thm:main}.
For general lattice coverings, Rogers obtained the bound
$n^{\log_2\log n+O(1)}$~\cite{Rog59}. Ordentlich, Regev, and Weiss
replaced it by the universal quadratic bound
$\theta_L(K)=O(n^2)$~\cite[Theorem~1.1]{ORW}. Li and Liu subsequently proved
\[
 \theta_L(K)\le Cn\log n\,(\log\log n)^{10/3+o(1)}
\]
uniformly over convex bodies~\cite[Theorem~1.2]{LL}.
Their binary-correction companion improves the exponent of
$\log\log n$ to
$\frac23\log_2(2\pi e)+o(1)$~\cite[Theorem~1.1]{LLC}.
Theorem~\ref{thm:main} removes this remaining iterated-logarithmic loss
and gives the constant-factor $n\log n$ bound.
In every sufficiently large dimension, Li and Liu also give centrally
symmetric bodies with lattice covering density at least
$cn\log n$~\cite[Theorem~1.3]{LL}. Together these
results determine the supremum of lattice covering density over bodies
in dimension $n$, up to absolute constant factors, as $n\log n$.
This statement concerns the supremum over bodies, rather than the covering
density of every individual body. A translative lower bound must exclude
all center configurations and is a separate question; the present theorem
and its proof concern the single-lattice upper bound.

\paragraph{The new mechanism.}
The proof develops the horizontal--vertical framework and binary-correction
ideas of Li and Liu~\cite{LL,LLC}. Write a point as $(x,y)$, where $x$
is the horizontal coordinate and the vertical coordinate $y$ has much
smaller dimension. Translating vertically lets several horizontal sections
of the same body contribute to covering a given height. The pointwise
Gaussian marginal theorem of Eldan and Klartag~\cite{EK} controls the
volumes of these sections after an affine normalization. The lattice
mean-hole method of Rogers and Schmidt~\cite{Rog58,Sch58}, in the
form of Ordentlich--Regev--Weiss~\cite[Theorem~2.3]{ORW}, then provides
one horizontal lattice with an exponentially small uncovered fraction for
every section on a finite list.

The remaining task is to distribute these horizontal contributions by one
linear shear of a vertical lattice. A binary correction offers two adjacent
integer choices in each vertical coordinate. In Li and Liu's two-stage
construction, one auxiliary correction is chosen for an entire weighted
Boolean cube~\cite[Theorem~6.2]{LLC}. Here the corrections themselves
form a hierarchy: we retain their weighted branches through all
intermediate blocks.

The total Gaussian mass of the retained family, multiplied by the
vertical lattice covolume, stays above an absolute constant,
while every individual atom becomes very small. These two properties
play different roles: the total mass supplies enough horizontal load,
and the atom bound keeps every section within the range of the mean-hole
estimate. The corrections are arranged in successively smaller blocks,
with sizes approximately $b,b^{0.9},b^{0.9^2},\ldots$, stopped at a
fixed absolute cutoff. Each block loses only $O(b^{-1/100})$ of the
mass. The resulting losses are summable even when the number of blocks
grows.

Two auxiliary statements isolate features of potential independent use.
The finite-group sampler in Lemma~\ref{lem:sampler} converts an
almost-everywhere average into a uniform lower average using random
subset sums; its cost depends on $\log|H|$, rather than $|H|$.
Lemma~\ref{lem:shear} extends the simultaneous Boolean-shift argument
of Li and Liu~\cite[Theorem~4.3]{LL} to finitely many binary-suffix
patterns with one linear shear. The values of the branch bases can depend
on later choices; only the unit separation of two siblings is needed.
This preserves the single-lattice structure throughout.

\paragraph{Organization.}
Section~\ref{sec:preliminaries} derives the precise section and horizontal
lattice estimates from two established results, stating their hypotheses
explicitly. Section~\ref{sec:sampler} proves the uniform finite-group
sampler. Sections~\ref{sec:blocks} and~\ref{sec:hierarchy} construct
the folded Gaussian blocks and their weighted branching hierarchy.
Section~\ref{sec:final} chooses the common shear and removes the final
holes by a dilation of factor $1+1/n$. All new probabilistic and
geometric arguments are proved below; no estimate on a single successful
branch substitutes for a bound on their total weight.

\paragraph{Notation and constants.}
Euclidean volume is written $\vol_d$. Expectations over finite groups
use uniform probability; $|A|$ denotes cardinality for a finite set.
On a lattice torus, $|A|$ instead denotes normalized Haar measure.
The standard Gaussian density is
\[
 \gamma_d(y)=(2\pi)^{-d/2}\exp(-\|y\|^2/2).
\]
All logarithms are natural except $\log_2$. The constants $c,C$ may
change from line to line and are absolute. Asymptotic statements concern
sufficiently large dimensions or block sizes; their thresholds never
depend on the convex body. The particular decimal exponents below are
chosen for strict inequalities, not optimization. A fixed cutoff may be
large, but it is chosen before the dimension and remains absolute.
